\documentclass[12pt]{article}
\usepackage[margin=1in]{geometry}
\usepackage{fancyhdr}
\usepackage{datetime}
\usepackage{amsmath}
\usepackage{booktabs}
\usepackage{siunitx}
\usepackage{tabularray}
\usepackage{float}
\usepackage{ulem}

% Header setup
\setlength{\headheight}{41.55003pt}
\addtolength{\topmargin}{-10pt}
\setlength{\headsep}{20pt}
\pagestyle{fancy}
\fancyhf{}
\renewcommand{\headrulewidth}{0pt}
\rhead{Brent Myers \\ Econ 5253 \\ 04/07/2026}

\begin{document}

\begin{center}
\textbf{\uline{Problem Set 8 - Matrix Manipulation, Optimization, and Simulated Data}}\\[0.5em]
\end{center}

\section*{Q4: Data Simulation}
I created a simulated dataset with $N = 100{,}000$ observations and $K = 10$ variables (including an intercept column of ones). The remaining 9 columns of $X$ were drawn from a standard normal distribution. The error term $\varepsilon$ was drawn from $N(0, 0.25)$ (i.e., $\sigma = 0.5$). The true parameter vector is:
\[
\beta = \begin{bmatrix} 1.5 & -1 & -0.25 & 0.75 & 3.5 & -2 & 0.5 & 1 & 1.25 & 2 \end{bmatrix}'
\]
The outcome variable was generated as $Y = X\beta + \varepsilon$.

\section*{Q5: OLS Closed-Form Solution}
Using the formula $\hat{\beta}_{OLS} = (X'X)^{-1}X'Y$, the estimated coefficients are extremely close to the true $\beta$. For example, $\hat{\beta}_1 = 1.501$ versus the true value of 1.5. With 100,000 observations, we expect very precise estimates.

\section*{Q6: OLS via Gradient Descent}
Using gradient descent with a learning rate of $\alpha = 0.0000003$, the algorithm converged at iteration 451. The resulting estimates are identical to the closed-form OLS solution, confirming that gradient descent successfully minimized the sum of squared residuals.

\section*{Q7: OLS via \texttt{nloptr} (L-BFGS and Nelder-Mead)}
Both the L-BFGS (gradient-based) and Nelder-Mead (derivative-free) algorithms produced estimates identical to the closed-form OLS solution. The answers do not differ between the two methods, though L-BFGS converges faster since it uses gradient information.

\section*{Q8: MLE via \texttt{nloptr} (L-BFGS)}
The MLE estimates of $\beta$ are identical to the OLS estimates, which is expected since OLS and MLE produce the same coefficient estimates under the assumption of normally distributed errors. The MLE estimate of $\sigma$ is 0.4999, which is very close to the true value of 0.5.

\section*{Q9: OLS via \texttt{lm()}}
Using R's built-in \texttt{lm()} function confirms all previous results. The regression table is shown below, followed by a comparison table of $\hat{\beta}$ estimates across all six methods.

\input{PS8_regression.tex}

All estimates of $\hat{\beta}$ are within a few thousandths of the ground truth $\beta$ used to generate the data, which is what we would expect given the large sample size ($N = 100{,}000$) and relatively small error variance ($\sigma^2 = 0.25$). As the comparison table below shows, all six estimation methods --- closed-form OLS, gradient descent, L-BFGS, Nelder-Mead, MLE, and \texttt{lm()} --- produce identical estimates to seven decimal places.

\begin{table}[H]
\centering
\begin{tblr}[         %% tabularray outer open
]                     %% tabularray outer close
{                     %% tabularray inner open
colspec={Q[]Q[]Q[]Q[]Q[]Q[]Q[]Q[]},
hline{2}={1-8}{solid, black, 0.05em},
hline{1}={1-8}{solid, black, 0.08em},
hline{12}={1-8}{solid, black, 0.08em},
}                     %% tabularray inner close
Beta & True & ClosedForm & GradDesc & LBFGS & NelderMead & MLE & lm \\
beta 1 & 1.5 & 1.5010518 & 1.5010518 & 1.5010518 & 1.5010518 & 1.5010518 & 1.5010518 \\
beta 2 & -1 & -1.0008296 & -1.0008296 & -1.0008296 & -1.0008296 & -1.0008296 & -1.0008296 \\
beta 3 & -0.25 & -0.251648 & -0.251648 & -0.251648 & -0.251648 & -0.251648 & -0.251648 \\
beta 4 & 0.75 & 0.7490406 & 0.7490406 & 0.7490406 & 0.7490406 & 0.7490406 & 0.7490406 \\
beta 5 & 3.5 & 3.5005531 & 3.5005531 & 3.5005531 & 3.5005531 & 3.5005531 & 3.5005531 \\
beta 6 & -2 & -2.0008185 & -2.0008185 & -2.0008185 & -2.0008185 & -2.0008185 & -2.0008185 \\
beta 7 & 0.5 & 0.4987148 & 0.4987148 & 0.4987148 & 0.4987148 & 0.4987148 & 0.4987148 \\
beta 8 & 1 & 1.0028269 & 1.0028269 & 1.0028269 & 1.0028269 & 1.0028269 & 1.0028269 \\
beta 9 & 1.25 & 1.2465102 & 1.2465102 & 1.2465102 & 1.2465102 & 1.2465102 & 1.2465102 \\
beta 10 & 2 & 2.0010012 & 2.0010012 & 2.0010012 & 2.0010012 & 2.0010012 & 2.0010012 \\
\end{tblr}
\end{table}


\end{document}